\cotefile{22}
\scanpage{2}
% printed running head
\textbf{C. R. Acad. Sc. Paris, t. 257, p. 2948--2950 (13 novembre 1963).}\hfill\textbf{Groupe 1.}\par\noindent\rule{\linewidth}{0.4pt}

\textsc{Algèbre.} --- \textit{Étude différentielle des anneaux locaux réguliers : applications.} Note (\textsuperscript{*}) de M. \textbf{\textsc{Jean-Pierre Jouanolou}}, présentée par M. René Garnier.

Le paragraphe 1 donne pour les anneaux locaux usuels (localités, localités analytiques, etc.) une caractérisation différentielle de régularité. Remarquant que, pour ces anneaux, la complétion se fait de façon « séparable », on montre que la réduction et la normalité analytiques en résultent simplement.

1. Les critères différentiels de régularité de Nakai, Suzuki et Mount font \deleted{apparaître}\inserted{\handwritten{appel}} en caractéristique $p>0$ \inserted{\handwritten{$/$ à}} une \deleted{condition}\inserted{\handwritten{$\angle$ hypothèse}} de réduction analytique (il y a des exemples contraires sinon). L’énoncé suivant permet d’éviter cette condition dans les cas usuels :
% The two oblique proof marks at the right repeat the insertion and substitution signs.

\margin{\handwritten{$\left.\rule{0pt}{7em}\right]$}}\textsc{Proposition 1.} --- \textit{Soient $\mathrm A$ un anneau local régulier équicaractéristique tel que \tikz[baseline=(x.base)]{\node[draw,ellipse,inner sep=1pt] (x) {$\mathrm D(\mathrm A)$};} soit plat et $\mathfrak a$ un idéal de $\mathrm A$ tel $\mathrm A/\mathfrak a$ soit \tikz[baseline=(x.base)]{\node[draw,ellipse,inner sep=1pt] (x) {\itshape réduit.};}}

\textit{Les assertions suivantes sont équivalentes :}

\textit{(a) $\mathrm A/\mathfrak a$ est régulier;}

\textit{(b) $\operatorname{Tor}^{1}_{\mathrm A/\mathfrak a}(\mathrm A/\mathfrak m,\ \mathrm D(\handwritten{\not}\mathrm A))=0$;}
\margin{\handwritten{$A/\underline{a}\,/$}}
% Matching oblique reference strokes touch the A inside D(A) in (b) and (c).

\textit{(c) $\operatorname{Tor}^{i}_{\mathrm A/\mathfrak a}(\mathrm A/\mathfrak m,\ \mathrm D(\handwritten{\not}\mathrm A))=0$ pour tout $i$ supérieur ou égal à $1$ ($\mathfrak m$ désignant l’idéal maximal de $\mathrm A$).}

\textit{De plus, si $\mathrm D(\mathrm A)$ est libre, (a) équivaut à :}

\textit{(d) $\mathrm D(\mathrm A/\mathfrak a)$ est libre.}

\textsc{Corollaire 1.} --- \textit{Soient $k$ un corps valué complet de caractéristique $p>0$ et tel que $(k:k^{p})<+\infty$, $\mathrm B$ une algèbre de séries convergentes sur $k$ et $\mathrm A$ un anneau réduit localisé d’un anneau quotient de $\mathrm B$. Pour que $\mathrm A$ soit régulier, il faut et il suffit que $\mathrm D(\mathrm A)$ soit libre.}

\textsc{Corollaire 2} [Kunz (\textsuperscript{3})]. --- \textit{Pour qu’une localité \uline{réduite} $\mathrm A$ soit régulière, il faut et il suffit que son module de différentielles absolu $\mathrm D(\mathrm A)$ soit libre.}
\margin{\handwritten{$\big|$ contenu dans le critère jacobien \ldots}}

\textsc{Corollaire 3} [Suzuki (\textsuperscript{10}), Mount (\textsuperscript{5})]. --- \textit{Soit $\mathrm A$ un anneau local nœthérien analytiquement \uline{réduit} de caractéristique $p>0$ et tel que le séparé $\mathrm D'(\mathrm A)$ soit de \uline{type fini}. Alors, pour que $\mathrm D'(\mathrm A)$ soit libre, il faut et il suffit que $\mathrm A$ soit régulier.}
\margin{\handwritten{assez rare !}}

La proposition 1 fournit des critères pour que l’ensemble des points non réguliers du spectre d’un anneau (lieu singulier) soit un fermé, notamment :

\margin{\handwritten{$\left.\rule{0pt}{2em}\right]$}}\textsc{Proposition 2.} --- \textit{Si $\mathrm A$ est un anneau global régulier équicaractéristique tel que $\mathrm D(\mathrm A)$ soit libre, le lieu singulier de toute $\mathrm A$-algèbre affine $\mathrm B$ est fermé.}
\margin{\handwritten{\unsure{séries|série} \unsure{formelles|formelle} ??}}

En particulier :

\textsc{Proposition 3.} --- \textit{Soit $\mathrm G$ un schéma en groupes régulier sur un corps parfait. Alors tout schéma de type fini sur $\mathrm G$ a un lieu singulier fermé.}
\margin{\handwritten{\unsure{canularesque|canularique|canulariques}}}

2. Dans tout ce qui suit, $\mathrm A$ étant un anneau local, nous désignerons par $\mathrm L_{\mathrm A}$ son corps résiduel et par $\mathrm K_{\mathrm A}$ son corps de fractions s’il est intègre.
% archive sheet number: 1
